\begin{frame}{Faster R-CNN: Ablauf im Projekt}
\small
\begin{columns}[T,onlytextwidth]
\begin{column}{0.46\textwidth}
\begin{enumerate}
    \item Backbone extrahiert Bildmerkmale.
    \item RPN schlägt mögliche Gesichtsregionen vor.
    \item ROI Head klassifiziert Regionen und verfeinert Boxen.
    \item IoU und Confidence entscheiden über Treffer.
\end{enumerate}
\end{column}
\begin{column}{0.50\textwidth}
\begin{itemize}
    \item WIDER FACE wurde in passende Annotationen für Torchvision übertragen.
    \item Im Red20-Smoke-Test lief Faster R-CNN gegen YOLOv8m, RetinaNet und FCOS.
    \item Ergebnis: guter Recall, aber höhere Latenz als schnellere One-Stage-Ansätze.
\end{itemize}
\end{column}
\end{columns}
\vspace{0.5em}
\scriptsize
\textbf{Unterschied:} Faster R-CNN sucht erst Kandidatenregionen und prüft sie danach genauer. YOLO/RetinaNet lokalisieren direkter; FCOS arbeitet zusätzlich anchor-free.
\end{frame}
